\section{Llegadas bajo la ley original}
Fijamos $\delta>0$, $M=\lceil\delta n\rceil$, $K=\lceil\sqrt n\rceil$ y
$S_M$ la entrada histórica regular. Desde la historia revelada hasta $M$ se sigue
$\nu$, con el orden futuro integrado, y se mata en la primera salida inclusiva.
La fila siguiente es uniforme. Usamos como entrada en cada padre regular la transición
en todos los hijos legales de B--C, dada por $p_{rc}(1+O(1/k))$ al fijar la fila,
con $p$ bistocástico y $\max p_{rc}\le C/k$. Los predictores hijos son finitos.
Cada fila tiene al menos una casilla de peso positivo incluso si $Q=0$. La existencia
de esta transición antes de parar viene de la regularidad admisible de padre e hijos.
Se presupone la brecha y difusión de torres en todas estas llegadas, como en R10, parte 04.

Dados $s$ índices de fila y $t$ de columna vivos en $M$ y un tamaño $j\ge K$,
escribimos $H_j^0$ para la llegada sin salida hasta $j$, incluido $S_M$.
Si algún índice falta ya en $M$, la prescripción tiene masa cero.
En un paso de tamaño $k>j$, las probabilidades de tocar alguna fila marcada o
alguna columna marcada son $s/k$ y al menos $t/k-Ct/k^2$, respectivamente.
La intersección cuesta a lo sumo $Cst/k^2$. Para $s,t\le j/2$,
el producto $(1-s/k)(1-t/k)$ es al menos $1/4$. La probabilidad de evitar
las marcas y seguir vivo se acota por ese producto por $\exp(C(st+t)/k^2)$.
Iterando en historias regulares y usando los factoriales descendentes,
\[
 \Pr_\nu(H^0_j,\ {\rm marcas\ vivas}\mid\mathcal F_M)
 \le (j/M)^{s+t}\exp(C(st+t)/j) \qquad\text{en }S_M .
\]
Se usó $(j)_s/(M)_s\le(j/M)^s$, y análogamente para las columnas.
La estimación no divide por la probabilidad de llegar ni por la de completar.

\paragraph{La excepción de una línea larga.}
En cada diagonal hay a lo sumo $M$ sitios candidatos ya en la entrada.
Para $R=\lceil\alpha j\rceil$, $m=\lfloor\alpha j/4\rfloor\ge1$,
cada selección de $m$ sitios exige $m$ filas y $m$ columnas distintas.
Comparando los momentos factoriales con $(R)_m$ y usando menos de $4n$ líneas,
\[
 \Pr_\nu(H_j^0,L_{\max}\ge R\mid\mathcal F_M)
 \le 4n\bigl(2j/(\alpha M)\bigr)^m\exp(C(m^2+m)/j).
\]
Fijemos el coste exterior $B_*\ge0$ de C, de modo que
$X=Qe^{-L}\le \exp(C_0+B_*j)$ en llegadas regulares, con $X=0$ si $Q=0$.
Se elige $b>0$ una sola vez, después de esas constantes y de $\alpha$, de modo que
\[
 x_0=2b/(\alpha\delta)<1,\qquad
 (\alpha/4)\log x_0+C\alpha^2/16<-(B_*+3).
\]
Para $K\le j\le\lfloor bn\rfloor$, se usa $m\log x_0\le(\alpha j/4-1)\log x_0$;
$\log(4n)/j=o(1)$ uniformemente. La probabilidad anterior es
$\le\exp(-(B_*+2)j)$ para $n$ grande. La parte 01 incluye toda llegada regular
sin completación en este evento. En particular también se obtiene esta cota
para tal llegada en $K$, bajo $\nu$, sin multiplicar por un índice de tiempos.
La esperanza de $1+X$ en la excepción completa es $O(e^{-2j})$.

\paragraph{Comparación con el argumento anterior.}
En la parada histórica hasta $K$ la ley inclinada coincide con $\nu$ antes de $M$
y tiene cociente de densidad entre $e^{-C/K}$ y $e^{C/K}$ después, incluso
en los hijos de salida (R9, partes 01--03). Se puede transferir una probabilidad
por ese factor una sola vez. El control anterior de la excepción se obtuvo bajo
$\nu$ directamente. No se afirma la cota para una historia que salió y reingresó.
